% !PRINTED 425
\weberpage{425}

Nach der Definition (11) ist $Q$ eine ganze rationale Function
von $r$, die, wenn man sie nach Potenzen von $r$ ordnet, ganze
rationale Zahlenco\"efficienten erh\"alt. Setzen wir also
\[
r = r_0^k,
\]
so wird
\[
Q = F(r_0^k),
\]
worin $F$ das Zeichen f\"ur eine rationale Function ist, deren
Co\"efficienten von $k$ nicht abh\"angig sind. Setzen wir nun f\"ur
$k$ der Reihe nach die Werthe
\[
k = 1,\ 2 \ldots n-1,
\]
so stellt $r_0^k$ alle Wurzeln der Kreistheilungsgleichung $X_n = 0$
dar, und die Summe
\[
\overset{k}{\Sigma} Q = F(r_0) + F(r_0^2) + \cdots + F(r_0^{n-1})
\]
ist eine \gesperrt{symmetrische Function} der Wurzeln dieser Gleichung; sie
l\"asst sich also rational durch die Co\"efficienten der Gleichung,
d.\,h. durch rationale Zahlen ausdr\"ucken und ist mithin selbst eine
\gesperrt{rationale Zahl}. Andererseits ist aber nach der Formel (15)
\[
\overset{k}{\Sigma} Q = h\, i^{\frac{n-1}{2}} \sqrt{n},
\]
wenn $h$ eine ganze Zahl bedeutet, n\"amlich die Anzahl der F\"alle,
in denen in (15) das positive Zeichen zu nehmen ist, vermindert um
die Anzahl der F\"alle, in denen das negative Zeichen gilt. Beides
ist aber nur dann mit einander vertr\"aglich, wenn $h = 0$ ist; und
damit ist der folgende Satz bewiesen:

\gesperrt{1. Durchl\"auft $k$ die Reihe der Zahlen 1, 2, 3, \dots\ $n-1$,
so gilt in der Formel (16) ebenso oft das positive wie das negative
Zeichen.}

Es braucht kaum besonders erw\"ahnt zu werden, dass man f\"ur $k$
auch ein anderes volles Restsystem von $n$, mit Ausschluss der durch
$n$ theilbaren Zahl, nehmen kann.

\begin{center}
% [corrected 2026-08-07, 10x zoom strips 8/10 and 9/10] The page prints "\S. 136." here. The
% fragment had written 137, a SILENT CORRECTION of what looked like a printer's slip -- exactly
% the class the spec forbids without a % [sic].
%
% And the correction was itself wrong. Section headings across the volume run
% ... 133 (i0439), 134 (i0442), [135], 136, 137 (i0457), 138 (i0464), 139 (i0474) ...
% With this fragment reading 137, the corpus carried TWO \S. 137 headings and no \S. 136 at all.
% The page is right and the transcription was wrong; restored to 136, no erratum needed.
%
% Method note: what settled this was not the crop alone but the crop plus the volume's own
% numbering sequence. A single heading cannot be judged in isolation -- a plausible repair looks
% exactly like a correct reading until you count the series.
\S.\,136.\\[2pt]
\gesperrt{Primitive Congruenzwurzeln.}
\end{center}

Parallel mit der Theorie der Einheitswurzeln geht eine Theorie der
sogenannten binomischen Congruenzen, ohne die in der Theorie der
Kreistheilungsgleichungen weitere Schritte nicht
%CONT>>
